\documentclass[12pt,a4paper]{ctexart}
\usepackage[top=2.5cm,bottom=2.5cm,left=2.5cm,right=2.5cm]{geometry}
\usepackage{amsmath,amssymb,amsthm}
\usepackage{physics}
\usepackage{graphicx}
\usepackage{float}
\usepackage{booktabs}
\usepackage{array}
\usepackage{caption}
\usepackage{hyperref}
\usepackage[numbers,sort&compress]{natbib}
\usepackage{xcolor}
\usepackage{enumitem}
\hypersetup{colorlinks=true,linkcolor=blue,citecolor=blue,urlcolor=blue}
\theoremstyle{definition}
\newtheorem{definition}{定义}[section]
\newtheorem{remark}{注记}[section]

\title{\heiti 格点QCD中的部分子分布函数（PDFs）：\\
从Feynman部分子模型到连续极限第一性原理计算}
\author{基于本库文献的综合解析}
\date{\today}

\begin{document}
\maketitle

\begin{abstract}
\noindent
部分子分布函数（PDFs）是描述强子内部夸克和胶子动量分布的基本物理量，是连接QCD第一性原理与高能对撞机实验唯象学的核心纽带。传统的PDF确定依赖于对深度非弹性散射（DIS）、Drell-Yan过程、喷注产生等大量实验数据的全局拟合（CT、NNPDF、MSHT等合作组），但这种方法在数据稀疏的大$x$区域面临本质性的模型依赖。格点QCD在大动量有效理论（LaMET）框架下的突破使得从第一性原理直接计算Bjorken-$x$依赖的PDF成为现实。本文系统梳理了PDF从连续定义到格点计算的完整知识链：涵盖光锥PDF的Feynman部分子模型起源与QCD因子化定义、全局拟合的现状与局限、LaMET框架下准PDF和赝PDF两种互补方案、格点上夸克和胶子PDF计算的完整技术链路（算符构造→关联函数→重整化→匹配→连续极限外推）、以及当前格点PDF计算与全局拟合结果的定量比较。重点展示了从单一格点间距的探索性计算到多格点间距连续极限外推的演进历程，并展望了在物理$\pi$介子质量下实现百分级精度的前景。
\end{abstract}

\tableofcontents
\newpage

\section{引言：PDF——连接QCD与实验的桥梁}

\subsection{Feynman的部分子模型}

部分子的概念源于Feynman在1969年提出的朴素部分子模型\cite{Ji2014LaMET}：一个以接近光速运动的强子（如质子）可以被视为一束互不相互作用的``部分子''——夸克和胶子——它们由动量密度$q(x)$和$g(x)$来表征，其中$x = k^z/P^z$是部分子携带的纵向动量份额。强子-强子硬散射截面可以表示为基本部分子散射截面$\hat{\sigma}$与部分子密度的卷积。

在QCD中，这一简单图像通过\textbf{因子化定理}得到了严格的场论论证\cite{Izubuchi2018}。唯一的修正来自量子场论：部分子密度是\textbf{方案和标度依赖的}，其标度演化由DGLAP（Dokshitzer-Gribov-Lipatov-Altarelli-Parisi）方程控制。

\subsection{PDF的光锥定义}

在现代QCD因子化框架中，非极化夸克PDF由光锥关联算符的核子矩阵元定义\cite{Izubuchi2018,Chen2025gluon}：

\begin{definition}[夸克PDF的光锥定义]
\begin{equation}
\boxed{q(x, \mu) = \int \frac{d\xi^-}{4\pi} e^{-ixP^+\xi^-} \langle P| \bar{\psi}(\xi^-) \gamma^+ \mathcal{U}(\xi^-, 0) \psi(0)(\mu) |P\rangle},
\label{eq:lightcone_quark_pdf}
\end{equation}
其中$\xi^\pm = (t \pm z)/\sqrt{2}$是光锥坐标，$P^+ = (P^0 + P^z)/\sqrt{2}$，$\mathcal{U}(\xi^-, 0) = \mathcal{P}\exp[-ig\int_0^{\xi^-} d\eta^- A^+(\eta^-)]$是光锥Wilson线。
\end{definition}

类似地，非极化胶子PDF由规范场张量的光锥关联定义\cite{Chen2025gluon,Balitsky2020}：

\begin{definition}[胶子PDF的光锥定义]
\begin{equation}
\boxed{xg(x, \mu) = \int \frac{dz^-}{2\pi P^+} e^{-ixP^+z^-} \langle P| F^{a}_{+\mu}(z^-) \mathcal{U}^{ab}(z^-, 0) F^{b\mu}_{\phantom{b\mu}+}(0)(\mu) |P\rangle},
\label{eq:lightcone_gluon_pdf}
\end{equation}
其中$\mathcal{U}^{ab}$是伴随表示中的光锥Wilson线。
\end{definition}

PDF满足\textbf{动量求和规则}（momentum sum rule）\cite{Fan2018gluon}：

\begin{equation}
\int_0^1 dx\, x\left[g(x,\mu) + \sum_{q=u,\bar{u},d,\ldots} q(x,\mu)\right] = 1,
\label{eq:momentum_sum_rule}
\end{equation}

即夸克和胶子携带的动量份额之和必须等于强子的总动量。这一求和规则是QCD中能量-动量守恒的直接体现。

\subsection{为什么格点QCD长期以来无法直接计算PDF？}

PDF的光锥定义涉及沿$\xi^- = (t-z)/\sqrt{2}$方向分离的场，\textbf{不可避地包含实时间依赖性}。格点QCD在Euclidean时空中表述（虚时间$\tau = it$），无法直接处理实时间关联。传统方法只能计算PDF的\textbf{Mellin矩}：

\begin{equation}
\int_0^1 dx\, x^{n-1} q(x, \mu) = a_n(\mu),
\label{eq:mellin_moments}
\end{equation}

其中$a_n(\mu)$是局域twist-2算符的约化矩阵元。但由于格点离散化破坏了$O(4)$旋转对称性，高维算符之间存在复杂的混合，实际只能可靠计算前两三个矩——远不足以重建完整的$x$依赖PDF。

\section{全局拟合：PDF的唯象学确定}

\subsection{方法论概述}

在LaMET方法成熟之前，PDF的主要来源是\textbf{全局拟合}（global fits）——即利用大量实验数据通过QCD因子化定理提取PDF。主要合作组包括：

\begin{table}[H]
\centering
\caption{主要PDF全局拟合合作组概览}
\label{tab:global_fits}
\begin{tabular}{lcccl}
\toprule
合作组 & 最新版本 & 微扰阶数 & 味数 & 方法 \\
\midrule
CT(EQ) & CT18 \cite{CT18} & NNLO & $2+1+1$ & Hessian \\
NNPDF & NNPDF4.0 \cite{NNPDF40} & NNLO & $2+1+1$ & 神经网络+MC \\
MSHT & MSHT20 \cite{MSHT20} & NNLO & $2+1+1$ & Hessian \\
JAM & JAM24 \cite{JAM24} & NLO/NNLO & $2+1+1$ & Monte Carlo \\
\midrule
CT14 & CT14 \cite{CT14} & NNLO & $2+1+1$ & Hessian \\
NNPDF3.1 & NNPDF3.1 \cite{NNPDF31} & NNLO & $2+1+1$ & 神经网络+MC \\
\bottomrule
\end{tabular}
\end{table}

全局拟合的核心思想是：在给定的参数化形式$f(x, Q_0) = x^{a}(1-x)^{b} \times (\text{多项式})$下，通过DGLAP演化计算每个实验观测量的理论预言，并最小化理论预言与实验数据之间的$\chi^2$。NNPDF合作组\cite{NNPDF40}采用机器学习（神经网络）替代传统参数化，实现了无偏的PDF提取，并通过``闭包检验''（closure test）——用已知的底层PDF生成伪数据再重新拟合——验证了方法的可靠性。

\subsection{实验数据输入}

全局拟合整合的实验数据类型极其丰富\cite{NNPDF40,NNPDF31,CT18}：

\begin{itemize}
    \item \textbf{深度非弹性散射（DIS）：}HERA的$e^\pm p$中性流和带电流数据——主要约束夸克PDF和胶子PDF（通过标度违规）；
    \item \textbf{固定靶Drell-Yan：}约束$\bar{d}/\bar{u}$不对称性；
    \item \textbf{对撞机弱矢量玻色子产生：}LHC和Tevatron的$W^\pm$和$Z$产生——区分上型和下型夸克；
    \item \textbf{单举喷注和双喷注产生：}LHC的ATLAS和CMS——主要约束中大$x$胶子PDF；
    \item \textbf{顶夸克对产生：}$t\bar{t}$微分分布——约束大$x$胶子；
    \item \textbf{孤立光子产生：}约束胶子PDF；
    \item \textbf{粲夸克产生：}约束小$x$胶子和粲PDF。
\end{itemize}

NNPDF4.0\cite{NNPDF40}在NNPDF3.1的基础上新增了44个数据集（主要来自LHC Run II），实现了对胶子和轻夸克味分离的显著改进。

\subsection{全局拟合的局限性与格点QCD的互补性}

尽管全局拟合取得了巨大成功，其本质局限在于\cite{Chen2025gluon}：

\begin{enumerate}
    \item \textbf{大$x$区域的模型依赖：}$x > 0.6$区域实验数据极度稀疏，PDF的行为主要由参数化形式的外推决定，不同合作组的结果差异显著；
    \item \textbf{胶子PDF的大$x$不确定性：}胶子PDF在大$x$处的全局拟合约束非常弱——因为大$x$胶子主要通过喷注和顶夸克产生间接约束，而这些过程的理论不确定性较大；
    \item \textbf{参数化偏差：}即使使用灵活的神经网络，有限的数据集也无法唯一确定PDF的全部函数形式。
\end{enumerate}

\textbf{格点QCD的第一性原理计算恰恰在这些方面具有独特的互补优势}：大$x$区域正是LaMET框架下格点计算精度最高的区域（因为大$x$对应于小$zP^z$，信噪比最佳）。如NNPDF合作组所述\cite{NNPDF40}：

\begin{quote}
``The path to proton structure at one-percent accuracy requires synergy between experimental data, lattice QCD, and perturbative calculations.''
\end{quote}

\section{LaMET框架：从Euclidean格点到光锥PDF}

\subsection{核心思想}

季向东在2013年提出的LaMET框架\cite{Ji2013Euclidean,Ji2014LaMET}彻底改变了格点QCD计算PDF的范式。其核心思想是利用Lorentz boost将Euclidean空间关联函数映射到光锥关联函数。

\textbf{准PDF（quasi-PDF）}的定义\cite{Ji2013Euclidean,Liu2019isovector}：

\begin{definition}[准夸克PDF]
\begin{equation}
\boxed{\tilde{q}(x, P^z) = \int_{-\infty}^{\infty} \frac{dz}{4\pi} e^{ixP^z z} \langle P| \bar{\psi}(z) \Gamma U(z, 0) \psi(0) |P\rangle},
\label{eq:quasi_pdf}
\end{equation}
其中$\Gamma = \gamma^t$（无算符混合）或$\gamma^z$，$U(z,0) = \mathcal{P}\exp[-ig\int_0^z dz' A^z(z')]$是空间方向的Wilson线。准PDF在大$P^z$下趋近光锥PDF。
\end{definition}

Liu等人\cite{Liu2019isovector}特别指出$\Gamma = \gamma^t$相对于$\Gamma = \gamma^z$的优势：$\gamma^t$算符在$O(a^0)$层次上不与标量准PDF算符$O_I$混合，从而避免了非微扰重整化中的额外系统误差。

\textbf{因子化定理}\cite{Izubuchi2018}将准PDF与光锥PDF联系起来：

\begin{equation}
\boxed{\tilde{q}(x, P^z, p^R_z, \mu_R) = \int_{-1}^{1} \frac{dy}{|y|} C\left(\frac{x}{y}, r, \frac{yP^z}{\mu}, \frac{yP^z}{p^R_z}\right) q(y, \mu) + \mathcal{O}\left(\frac{M^2}{P_z^2}, \frac{\Lambda^2_{\text{QCD}}}{P_z^2}\right)},
\label{eq:factorization}
\end{equation}

其中$C$是微扰可计算的匹配核。

\subsection{准PDF vs 赝PDF：两种互补方案}

\begin{table}[H]
\centering
\caption{准PDF（LaMET）与赝PDF方案的对比}
\label{tab:quasi_vs_pseudo}
\begin{tabular}{lcc}
\toprule
特性 & 准PDF（LaMET）& 赝PDF（pseudo-PDF）\\
\midrule
基本量 & $\tilde{q}(x, P^z)$（动量空间）& $\mathfrak{M}(\nu, z^2)$（坐标空间）\\
Ioffe时间 & $\nu = zP^z$ & $\nu = -(p \cdot z)$\\
大标度极限 & $P^z \to \infty$ & $z \to 0$（短距离）\\
因子化 & 动量空间卷积 & 坐标空间乘积\\
重整化 & RI/MOM或混合方案 & ratio方案（重整化因子抵消）\\
Fourier变换 & 坐标→动量 & 无需（直接分析$\nu$空间）\\
主要合作组 & LPC, LP$^3$, CLQCD, MSULat & HadStruc, MSULat\\
\bottomrule
\end{tabular}
\end{table}

Izubuchi等人\cite{Izubuchi2018}严格证明了\textbf{两种方案的因子化公式等价}——它们从相同的格点可观测量出发，只是在不同空间（动量vs坐标）中执行因子化。

\subsection{胶子PDF的LaMET计算}

对于胶子PDF，准PDF定义为\cite{Chen2025gluon,Fan2018gluon,Zhang2019}：

\begin{equation}
\tilde{x}\tilde{g}(x, P^z) = \frac{1}{\mathcal{N}} \int \frac{dz}{2\pi P^z} e^{-ixzP^z} \tilde{h}^R(z, P^z),
\label{eq:gluon_quasi_pdf}
\end{equation}

其中重整化矩阵元$\tilde{h}^R(z, P^z) = \langle P|\mathcal{O}(z)|P\rangle_R$由胶子算符（非极化或极化）的核子矩阵元给出。归一化因子$\mathcal{N} = (P^z)^2/(P^t)^2$。

在基础表示下，胶子算符取可乘性重整化形式\cite{Chen2025gluon,Zhang2019}：

\begin{equation}
\mathcal{O}(z) = M_{tx;tx}(z) + M_{ty;ty}(z) - 2M_{xy;xy}(z),
\label{eq:gluon_operator}
\end{equation}

其中$M_{\mu\lambda;\nu\rho}(z) = \sum_{\vec{x}} \operatorname{Tr}[F_{\mu\lambda}(\vec{x}+z\hat{z})U(\vec{x}+z\hat{z},\vec{x})F_{\nu\rho}(\vec{x})U(\vec{x},\vec{x}+z\hat{z})]$。

\textbf{胶子准PDF的关键挑战}包括：
\begin{itemize}
    \item 非连通图（胶子流与核子无夸克线连接）导致信噪比远差于夸克情形\cite{Chen2025gluon}；
    \item 不同Lorentz分量具有不同反常量纲，需仔细构造可乘性重整化组合\cite{Zhang2019}；
    \item 线性发散因$C_A/C_F$因子而比夸克情形更强\cite{Chen2025gluon}。
\end{itemize}

\section{格点PDF计算的完整技术链路}

\subsection{从组态生成到裸矩阵元}

完整的格点PDF计算始于\textbf{规范组态的生成}。CLQCD合作组\cite{CLQCD2024,CLQCD2025}使用TITLS（tadpole改进树级Symanzik）规范作用量和TITLC（tadpole改进树级Clover）费米子作用量生成了覆盖$a \in [0.05, 0.11]$~fm、$m_\pi \in [135, 350]$~MeV的广泛系综。MSULat组\cite{NieMiera2025,Lin2025gluon}使用MILC合作组的HISQ系综（$N_f = 2+1+1$）。

\textbf{夸克传播子}的计算是计算量最大的步骤。蒸馏方法\cite{Peardon2009}通过截断Laplace算符的本征模（$N_{\text{ev}} \sim 64$--$100$），在计算代价与信噪比之间实现平衡。高动量boost通过相位因子$\tilde{v}^k_x = e^{i\vec{\zeta}\cdot\vec{z}}v^k_x$实现\cite{Egerer2021a}。

\textbf{胶子流算符}的构造涉及Clover $F_{\mu\nu}$和Wilson线的离散化。HYP涂抹（典型5--10步）或梯度流（$\tau_W \sim 3a^2$）被用于抑制Wilson线的线性紫外发散\cite{Chen2025gluon,NieMiera2025}。

\subsection{重整化：从裸到重整化矩阵元}

裸矩阵元$\tilde{h}^B(z, P^z, 1/a)$包含线性发散（$\sim e^{-\delta m|z|}$）和对数紫外发散。LPC合作组的\textbf{混合重整化方案}\cite{Ji2021hybrid}在短距离使用ratio方案（除以零动量矩阵元），在长距离使用自重整化：

\begin{equation}
\tilde{h}^R(z, P^z) = \frac{\tilde{h}^n_B(z, P^z)}{\tilde{h}^n_B(z, 0)} \theta(z_s - |z|) + \eta_s \frac{\tilde{h}^n_B(z, P^z)}{Z_R(z, 1/a)} \theta(|z| - z_s)。
\label{eq:hybrid}
\end{equation}

自重整化因子$Z_R$的参数通过在多格点间距上拟合$P^z=0$的零动量数据确定\cite{Huo2021self,Chen2025gluon}。RI/MOM方案\cite{Liu2019isovector,Zhang2019}提供了另一种非微扰重整化途径。

\subsection{匹配：从准PDF到光锥PDF}

重整化矩阵元经Fourier变换得到准PDF后，通过微扰匹配核转换为$\overline{\text{MS}}$ PDF。匹配核的精度取决于：

\begin{itemize}
    \item \textbf{微扰阶数：}当前标准为NLO，NNLO正在发展中\cite{Yao2023}；
    \item \textbf{对数重求和：}RGR（重整化群重求和）控制大对数$\alpha_s^n\ln^n(P^z/\mu)$\cite{Zhang2023RGR}；
    \item \textbf{重正子重求和：}LRR（领头重正子重求和）控制$1/P^z$幂次修正中的非微扰模糊性\cite{Zhang2023RGR,Su2023LRR}。
\end{itemize}

对于胶子PDF，匹配涉及$2 \times 2$的味混合矩阵\cite{Zhang2019,Yao2023}：

\begin{equation}
\begin{pmatrix} \tilde{g}(x, P^z) \\ \tilde{q}(x, P^z) \end{pmatrix} = \int \frac{dy}{y}
\begin{pmatrix} C_{gg} & C_{gq} \\ C_{qg} & C_{qq} \end{pmatrix}
\begin{pmatrix} g(y, \mu) \\ q(y, \mu) \end{pmatrix} + \cdots。
\label{eq:mixing_matrix}
\end{equation}

\subsection{连续极限与无穷大动量外推}

经过匹配的PDF仍然包含格点间距$a$和核子动量$P^z$的残余依赖性。最终的物理PDF通过联合外推获得\cite{Chen2025gluon}：

\begin{equation}
\boxed{xg(x, P^z, a) = xg_0(x) + a^2 f(x) + a^2 (P^z)^2 h(x) + \frac{d(x)}{(P^z)^2}}。
\label{eq:joint_extrap}
\end{equation}

这一外推是当前格点PDF计算中\textbf{最大的系统误差来源}。需要至少三个格点间距和多个核子动量的数据才能可靠地执行。

\section{格点PDF计算的主要成果}

\subsection{夸克PDF}

\textbf{LP$^3$合作组}\cite{LP3helicity,Lin2018}在物理$\pi$介子质量下完成了质子同位旋矢量夸克螺旋度分布$\Delta u(x) - \Delta d(x)$的计算，核子动量高达3.0~GeV，结果与NNPDF和JAM的全局分析在大$x$区域$1\sigma$内一致。

\textbf{LPC合作组}\cite{Liu2019isovector}使用$\Gamma = \gamma^t$准PDF算符系统研究了非极化同位旋矢量夸克PDF，详细分析了RI/MOM标度（$\mu_R$, $p^R_z$）依赖性、匹配中的投影算符选择，以及核子动量$P^z$依赖性。该工作确认了$\gamma^t$算符相对$\gamma^z$算符在避免算符混合方面的优势。

\subsection{胶子PDF——从探路到连续极限}

胶子PDF的格点计算经历了快速演进：

\begin{enumerate}
    \item \textbf{Fan et al.\ (2018)} \cite{Fan2018gluon}：首次尝试，$a = 0.111$~fm（单一格点间距），$m_\pi = 340$~MeV，$P^z \leq 0.93$~GeV，ratio重整化。结果在统计误差内与CT14和PDF4LHC15一致，建立了可行性。

    \item \textbf{MSULat (2023--2025)} \cite{Fan2023,Good2025a,Lin2025gluon}：扩展到$2+1+1$味HISQ系综，梯度流涂抹，赝PDF和LaMET双方案。在物理$\pi$介子质量附近的初步结果显示大$x$区胶子密度趋于零。

    \item \textbf{CLQCD/LPC (2025)} \cite{Chen2025gluon,Chen2025first}：首次\textbf{三格点间距连续极限外推}，$a = \{0.105, 0.0897, 0.0775\}$~fm，$m_\pi \approx 300$~MeV，$P^z$最高1.97~GeV，蒸馏+混合重整化+HYP5涂抹。最终结果$xg(x)/\langle x \rangle_g$与CT18 NNLO、NNPDF NNLO和JAM24在误差范围内一致，大$x$区域胶子PDF被压低至零附近。

    \item \textbf{MSULat自重整化 (2025)} \cite{NieMiera2025}：三格点间距梯度流自重整化，$a = \{0.151, 0.121, 0.089\}$~fm，$P^z \approx 2.0$--$2.2$~GeV，$m_\pi \approx 310$~MeV。连续极限结果与选择近零大$x$胶子密度的全局拟合一致。
\end{enumerate}

\subsection{胶子PDF格点结果与全局拟合的比较}

Chen等人\cite{Chen2025gluon}的连续极限外推结果（原文图4）展示了格点QCD与全局拟合在非极化胶子PDF上的定量比较：

\begin{itemize}
    \item 在中等$x$（$0.2 \lesssim x \lesssim 0.6$）区域：格点结果与CT18、NNPDF、JAM24在误差范围内一致；
    \item 在大$x$（$x > 0.6$）区域：格点和全局拟合的中心值都接近零，表明胶子PDF在大$x$处被强烈压低；
    \item 格点误差在大$x$和小$x$两端迅速增大——这是LaMET方法的固有局限（大$x$受限于有限$a$，小$x$受限于有限$P^z$）。
\end{itemize}

\subsection{介子PDF}

除了核子PDF外，格点QCD也已成功应用于介子PDF的计算：

\begin{itemize}
    \item \textbf{$\pi$介子胶子PDF：}Fan等人\cite{Fan2023pion}和MSULat组\cite{Lin2025gluon}的计算显示$\pi$介子胶子动量份额在大$x$处比核子更大，与QCD演化方程的预期一致。
    \item \textbf{$K$介子胶子PDF：}首次探索性计算\cite{Lin2025kaon}正在开展。
    \item \textbf{氘子类双重子系统PDF：}LPC合作组\cite{Chu2024deuteron}的计算验证了LaMET框架对多强子系统的适用性。
\end{itemize}

\section{拓展：广义部分子分布与三维结构}

PDFs只是部分子物理的``一维''投影（仅依赖于纵向动量份额$x$）。完整的三维强子结构需要以下推广\cite{Ji2013Euclidean}：

\begin{itemize}
    \item \textbf{广义部分子分布（GPDs）：}依赖于$x$（纵向）、$\xi$（动量传递的斜度）和$t$（动量传递平方）。GPDs可以通过深度虚Compton散射（DVCS）实验测量，并在LaMET框架下从格点QCD计算\cite{Yao2023}。

    \item \textbf{横向动量依赖分布（TMDs）：}依赖于$x$和$\vec{k}_\perp$（横向动量）。TMDs通过半单举深度非弹性散射（SIDIS）和Drell-Yan过程测量。在格点上，TMDs可以通过包含空间方向Wilson线（``staple''）的算符来计算，Collins-Soper核的格点确定已取得重要进展\cite{Ji2013Euclidean}。

    \item \textbf{分布振幅（DAs）：}$\pi$介子和$K$介子的leading-twist光锥分布振幅已在LaMET框架下被成功计算\cite{Yao2023,Ji2013Euclidean}。

    \item \textbf{Wigner分布：}作为GPDs和TMDs的生成函数，Wigner分布提供了最完整的强子内部结构信息，但其格点计算仍处于概念阶段\cite{Ji2013Euclidean}。
\end{itemize}

\section{从现有精度到百分级精度：未来的路径}

NNPDF合作组的路线图\cite{NNPDF40}提出了将质子结构确定推进到百分级精度的宏伟目标。实现这一目标需要理论和实验的多方面进展，而格点QCD将在其中发挥不可替代的作用：

\begin{enumerate}
    \item \textbf{物理$\pi$介子质量计算：}当前大部分胶子PDF计算仍在$m_\pi \sim 300$~MeV下进行。在物理$m_\pi \approx 140$~MeV下的计算，结合手征外推，是减少系统误差的关键。

    \item \textbf{更细的格点间距和更高动量：}$a \lesssim 0.04$~fm、$P^z \gtrsim 3$~GeV的数据将显著改善连续极限和无穷大动量外推的可靠性。

    \item \textbf{NNLO匹配核：}对夸克和胶子准PDF算符的NNLO匹配核计算（类似于已完成的同位旋矢量非极化夸克情形）将减少微扰匹配中的标度不确定性。

    \item \textbf{系统化的重正子处理：}将RGR+LRR方案从$\pi$介子夸克PDF\cite{Zhang2023RGR}推广到核子胶子PDF。

    \item \textbf{格点数据与全局拟合的联合分析：}将格点QCD在大$x$区域的精确约束纳入全局拟合框架，实现实验数据与第一性原理计算的协同。

    \item \textbf{EIC时代的到来：}电子-离子对撞机将以前所未有的精度测量极化DIS，为胶子和海夸克PDF提供决定性约束，特别是小$x$区域。
\end{enumerate}

\begin{quote}
``The path to proton structure at one-percent accuracy represents a confluence of experimental precision, theoretical sophistication, and computational power that has never before been achieved in the study of strongly interacting matter.''\cite{NNPDF40}
\end{quote}

\section{总结}

部分子分布函数是连接QCD第一性原理与高能对撞机实验唯象学的核心物理量。本文系统梳理了PDF从Feynman的朴素部分子模型到现代格点QCD连续极限计算的完整演化历程。

\textbf{核心演进脉络}可以概括为三条并行推进的线索：

\begin{enumerate}
    \item \textbf{唯象学线索：}从早期CTEQ的Hessian拟合到NNPDF4.0的机器学习神经网络，全局拟合的精度和稳健性不断提升。然而大$x$胶子PDF的约束仍然薄弱，为格点QCD的介入提供了明确的需求。

    \item \textbf{方法论线索：}LaMET框架（2013年提出）→因子化定理证明（2018年）→混合重整化方案（2021年）→RGR+LRR领先幂次精度匹配（2023年）。每一次理论突破都显著提升了格点PDF计算的系统精度控制水平。

    \item \textbf{计算实践线索：}夸克同位旋矢量PDF（2018--2020年，物理$m_\pi$）→胶子准PDF首次尝试（2018年，$a=0.111$~fm）→胶子PDF自重整化（2021--2025年）→胶子PDF三格点间距连续极限（2025年）。计算规模和精度在短短七年内实现了跨越式发展。
\end{enumerate}

\textbf{当前状态：}格点QCD计算的胶子PDF已首次达到与全局拟合在定量上一致的水平，同时进行了连续极限外推。这是格点部分子物理领域的一个里程碑，标志着从``可行性验证''到``精度物理''的范式转换。然而，在物理$\pi$介子质量和大$P^z$动量下的系统计算仍需大量资源和理论发展。

\textbf{下一阶段的目标：}在物理$\pi$介子质量、多个格点间距和更高核子动量下，将胶子PDF的格点计算推进到可与最精确的全局拟合竞争的精度水平，并将格点结果系统地整合入全局拟合框架。这将最终实现核子一维和三维部分子结构从QCD第一性原理的完整确定。

\begin{thebibliography}{99}

% ===== 全局拟合 =====
\bibitem{CT18}
T.-J.~Hou \textit{et al.},
``New CTEQ global analysis of quantum chromodynamics with high-precision data from the LHC,''
Phys.\ Rev.\ D \textbf{103}, 014013 (2021), arXiv:1912.10053.
\quad\textbf{[来源:} \texttt{docs/New CTEQ global analysis of quantum chromodynamics with high-precision data from the LHC.pdf}\textbf{]}

\bibitem{NNPDF40}
R.~D.~Ball \textit{et al.} (NNPDF),
``The path to proton structure at 1\% accuracy,''
Eur.\ Phys.\ J.\ C \textbf{82}, 428 (2022), arXiv:2109.02653.
\quad\textbf{[来源:} \texttt{docs/The Path to Proton Structure at One-Percent Accuracy.pdf}\textbf{]}

\bibitem{NNPDF31}
R.~D.~Ball \textit{et al.} (NNPDF),
``Parton distributions from high-precision collider data,''
Eur.\ Phys.\ J.\ C \textbf{77}, 663 (2017), arXiv:1706.00428.
\quad\textbf{[来源:} \texttt{docs/Parton distributions from high-precision collider data.pdf}\textbf{]}

\bibitem{MSHT20}
S.~Bailey, T.~Cridge, L.~A.~Harland-Lang, A.~D.~Martin, and R.~S.~Thorne,
``Parton distributions from LHC, HERA, Tevatron and fixed target data: MSHT20 PDFs,''
Eur.\ Phys.\ J.\ C \textbf{81}, 341 (2021), arXiv:2012.04684.
\quad\textbf{[来源:} \texttt{docs/Parton distributions from LHC, HERA, Tevatron and fixed target data MSHT20 PDFs.pdf}\textbf{]}

\bibitem{CT14}
S.~Dulat \textit{et al.},
``New parton distribution functions from a global analysis of quantum chromodynamics,''
Phys.\ Rev.\ D \textbf{93}, 033006 (2016), arXiv:1506.07443.
\quad\textbf{[来源:} \texttt{docs/New parton distribution functions from a global analysis of quantum chromodynamics.pdf}\textbf{]}

\bibitem{JAM24}
T.~Anderson, W.~Melnitchouk, and N.~Sato,
``New global analysis of spin-dependent parton distribution functions,''
(2024), arXiv:2501.00665.

% ===== LaMET框架 =====
\bibitem{Ji2013Euclidean}
X.~Ji,
``Parton physics on a Euclidean lattice,''
Phys.\ Rev.\ Lett.\ \textbf{110}, 262002 (2013), arXiv:1305.1539.
\quad\textbf{[来源:} \texttt{docs/Parton Physics on Euclidean Lattice.pdf}\textbf{]}

\bibitem{Ji2014LaMET}
X.~Ji,
``Parton physics from large-momentum effective field theory,''
Sci.\ China Phys.\ Mech.\ Astron.\ \textbf{57}, 1407 (2014), arXiv:1404.6680.
\quad\textbf{[来源:} \texttt{docs/Parton Physics from Large-Momentum Effective Field Theory.pdf}\textbf{]}

\bibitem{Izubuchi2018}
T.~Izubuchi, X.~Ji, L.~Jin, I.~W.~Stewart, and Y.~Zhao,
``Factorization theorem relating Euclidean and light-cone parton distributions,''
Phys.\ Rev.\ D \textbf{98}, 056004 (2018), arXiv:1801.03917.
\quad\textbf{[来源:} \texttt{docs/Factorization Theorem Relating Euclidean and Light-Cone Parton Distributions.pdf}\textbf{]}

% ===== 夸克PDF格点计算 =====
\bibitem{Liu2019isovector}
Y.-S.~Liu \textit{et al.} (Lattice Parton Collaboration (LPC)),
``Unpolarized isovector quark distribution function from lattice QCD: A systematic analysis of renormalization and matching,''
Phys.\ Rev.\ D \textbf{101}, 034020 (2020), arXiv:1807.06566.
\quad\textbf{[来源:} \texttt{docs/Unpolarized isovector quark distribution function from Lattice QCD A systematic analysis of renormalization and matching.pdf}\textbf{]}

\bibitem{LP3helicity}
H.-W.~Lin \textit{et al.} (LP$^3$ Collaboration),
``Proton isovector helicity distribution on the lattice at physical pion mass,''
Phys.\ Rev.\ Lett.\ \textbf{121}, 242003 (2018), arXiv:1807.07431.
\quad\textbf{[来源:} \texttt{docs/Proton Isovector Helicity Distribution on the Lattice at Physical Pion Mass.pdf}\textbf{]}

\bibitem{Lin2018}
H.-W.~Lin \textit{et al.},
``Parton distributions and lattice QCD calculations: toward 3D structure,''
Prog.\ Part.\ Nucl.\ Phys.\ \textbf{100}, 107 (2018), arXiv:1711.07916.

% ===== 胶子PDF格点计算 =====
\bibitem{Chen2025gluon}
C.~Chen \textit{et al.} (CLQCD, Lattice Parton),
``Unpolarized gluon PDF of the nucleon from lattice QCD in the continuum limit,''
(2025), arXiv:2510.26425.
\quad\textbf{[来源:} \texttt{docs/Unpolarized gluon PDF of the nucleon from lattice QCD in the continuum limit.pdf}\textbf{]}

\bibitem{Chen2025first}
C.~Chen \textit{et al.} (CLQCD, Lattice Parton),
``First self-renormalized gluon PDF of nucleon from large-momentum effective theory in the continuum limit,''
Phys.\ Rev.\ D \textbf{111}, 074506 (2025), arXiv:2408.12819.
\quad\textbf{[来源:} \texttt{docs/First Self-Renormalized Gluon PDF of Nucleon from Large-Momentum Effective Theory in the Continuum Limit.pdf}\textbf{]}

\bibitem{NieMiera2025}
A.~NieMiera, W.~Good, H.-W.~Lin, and F.~Yao,
``First self-renormalized gluon PDF of nucleon from large-momentum effective theory in the continuum limit,''
(2025), arXiv:2510.17758.
\quad\textbf{[来源:} \texttt{docs/First Self-Renormalized Gluon PDF of Nucleon from Large-Momentum Effective Theory in the Continuum Limit.pdf}\textbf{]}

\bibitem{Fan2018gluon}
Z.-Y.~Fan, Y.-B.~Yang, A.~Anthony, H.-W.~Lin, and K.-F.~Liu,
``Gluon quasi-PDF from lattice QCD,''
Phys.\ Rev.\ Lett.\ \textbf{121}, 242001 (2018), arXiv:1808.02077.
\quad\textbf{[来源:} \texttt{docs/Gluon Quasi-PDF From Lattice QCD.pdf}\textbf{]}

\bibitem{Good2025a}
W.~Good, F.~Yao, and H.-W.~Lin,
``First nucleon gluon PDF from large momentum effective theory,''
(2025), arXiv:2505.13321.
\quad\textbf{[来源:} \texttt{docs/First Nucleon Gluon PDF from Large Momentum Effective Theory.pdf}\textbf{]}

\bibitem{Lin2025gluon}
W.~Good, K.~Hasan, and H.-W.~Lin,
``Gluon parton distribution of the nucleon from $2+1+1$-flavor lattice QCD in the physical-continuum limit,''
J.\ Phys.\ G \textbf{52}, 035105 (2025), arXiv:2409.02750.
\quad\textbf{[来源:} \texttt{docs/Gluon Parton Distribution of the Nucleon from 2+1+1-Flavor Lattice QCD in the Physical-Continuum Limit.pdf}\textbf{]}

\bibitem{Fan2023}
Z.~Fan, W.~Good, and H.-W.~Lin,
``Physical-continuum limit of the nucleon gluon parton distribution from lattice QCD,''
Phys.\ Rev.\ D \textbf{108}, 014508 (2023), arXiv:2210.09985.
\quad\textbf{[来源:} \texttt{docs/Physical-Continuum Limit of the Nucleon Gluon Parton Distribution from Lattice QCD.pdf}\textbf{]}

\bibitem{Fan2023pion}
Z.~Fan and H.-W.~Lin,
``Gluon parton distribution of the pion from lattice QCD,''
Phys.\ Lett.\ B \textbf{823}, 136778 (2021), arXiv:2104.06372.
\quad\textbf{[来源:} \texttt{docs/Gluon Parton Distribution of the Pion from Lattice QCD.pdf}\textbf{]}

% ===== 胶子算符 =====
\bibitem{Zhang2019}
J.-H.~Zhang, X.~Ji, A.~Sch\"afer, W.~Wang, and S.~Zhao,
``Accessing gluon parton distributions in large momentum effective theory,''
Phys.\ Rev.\ Lett.\ \textbf{122}, 142001 (2019), arXiv:1808.10824.
\quad\textbf{[来源:} \texttt{docs/Accessing gluon parton distributions in large momentum effective theory.pdf}\textbf{]}

\bibitem{Balitsky2020}
I.~Balitsky, W.~Morris, and A.~Radyushkin,
``Gluon pseudo-distributions at short distances: Forward case,''
Phys.\ Lett.\ B \textbf{808}, 135621 (2020), arXiv:1910.13963.
\quad\textbf{[来源:} \texttt{docs/Short-Distance Structure of Unpolarized Gluon Pseudodistributions.pdf}\textbf{]}

% ===== 重整化 =====
\bibitem{Ji2021hybrid}
X.~Ji, Y.~Liu, A.~Sch\"afer, W.~Wang, Y.-B.~Yang, J.-H.~Zhang, and Y.~Zhao,
``A hybrid renormalization scheme for quasi light-front correlations in large-momentum effective theory,''
Nucl.\ Phys.\ B \textbf{964}, 115311 (2021), arXiv:2008.03886.
\quad\textbf{[来源:} \texttt{docs/A Hybrid Renormalization Scheme for Quasi Light-Front Correlations in Large-Momentum Effective Theory.pdf}\textbf{]}

\bibitem{Huo2021self}
Y.-K.~Huo \textit{et al.} (LPC),
``Self-renormalization of quasi-light-front correlators on the lattice,''
Nucl.\ Phys.\ B \textbf{969}, 115443 (2021), arXiv:2103.02965.
\quad\textbf{[来源:} \texttt{docs/Self-Renormalization of Quasi-Light-Front Correlators on the Lattice.pdf}\textbf{]}

% ===== 精度控制 =====
\bibitem{Zhang2023RGR}
R.~Zhang, J.~Holligan, X.~Ji, and Y.~Su,
``Leading power accuracy in lattice calculations of parton distributions,''
Phys.\ Lett.\ B \textbf{844}, 138081 (2023), arXiv:2305.05212.
\quad\textbf{[来源:} \texttt{docs/Leading Power Accuracy in Lattice Calculations of Parton Distributions.pdf}\textbf{]}

\bibitem{Su2023LRR}
Y.~Su, J.~Holligan, X.~Ji, F.~Yao, J.-H.~Zhang, and R.~Zhang,
``Power corrections and renormalons in parton quasi-distributions,''
Nucl.\ Phys.\ B \textbf{991}, 116201 (2023), arXiv:2209.01236.
\quad\textbf{[来源:} \texttt{docs/Power corrections and renormalons in parton quasi-distributions.pdf}\textbf{]}

% ===== 匹配 =====
\bibitem{Yao2023}
F.~Yao, Y.~Ji, and J.-H.~Zhang,
``Connecting Euclidean to light-cone correlations,''
JHEP \textbf{11}, 021 (2023), arXiv:2212.14415.
\quad\textbf{[来源:} \texttt{docs/Connecting Euclidean to light-cone correlations From flavor nonsinglet in forward kinematics to flavor singlet in non-forward kinematics.pdf}\textbf{]}

% ===== 蒸馏 =====
\bibitem{Peardon2009}
M.~Peardon \textit{et al.} (Hadron Spectrum),
``A novel quark-field creation operator construction for hadronic physics in lattice QCD,''
Phys.\ Rev.\ D \textbf{80}, 054506 (2009), arXiv:0905.2160.
\quad\textbf{[来源:} \texttt{docs/A novel quark-field creation operator construction for hadronic physics in lattice QCD.pdf}\textbf{]}

\bibitem{Egerer2021a}
C.~Egerer, R.~G.~Edwards, K.~Orginos, and D.~G.~Richards,
``Distillation at high-momentum,''
Phys.\ Rev.\ D \textbf{103}, 034502 (2021), arXiv:2009.10691.
\quad\textbf{[来源:} \texttt{docs/Distillation at High-Momentum.pdf}\textbf{]}

% ===== 组态 =====
\bibitem{CLQCD2024}
Z.-C.~Hu \textit{et al.} (CLQCD),
``Charmed meson masses and decay constants in the continuum,''
Phys.\ Rev.\ D \textbf{109}, 054507 (2024), arXiv:2310.00814.
\quad\textbf{[来源:} \texttt{docs/Charmed meson masses and decay constants in the continuum from the tadpole improved clover ensembles.pdf}\textbf{]}

\bibitem{CLQCD2025}
H.-Y.~Du \textit{et al.} (CLQCD),
``Quark masses and low energy constants in the continuum,''
Phys.\ Rev.\ D \textbf{111}, 054504 (2025), arXiv:2408.03548.
\quad\textbf{[来源:} \texttt{docs/Quark masses and low energy constants in the continuum from the tadpole improved clover ensembles.pdf}\textbf{]}

% ===== 介子/多强子系统 =====
\bibitem{Lin2025kaon}
``First glimpse into the kaon gluon parton distribution using lattice QCD,''
\quad\textbf{[来源:} \texttt{docs/First Glimpse into the Kaon Gluon Parton Distribution Using Lattice QCD.pdf}\textbf{]}

\bibitem{Chu2024deuteron}
M.-H.~Chu \textit{et al.} (LPC),
``Parton distribution function of a deuteron-like dibaryon system from lattice QCD,''
\quad\textbf{[来源:} \texttt{docs/Parton Distribution Function of a Deuteron-like Dibaryon System from Lattice QCD.pdf}\textbf{]}

\end{thebibliography}

\end{document}
